% GENERATED FILE. Edit canonical lessons, then run npm run generate:books.
% canonical-source-hash: fnv1a64:bad4ec1e35ef7ce9
% canonical-lessons: ES-C648-parquimetro, ES-C648-estacionamiento, ES-C648-remolque, ES-C648-carril, ES-C648-remolcar

\chapter{Dónde dejar el coche}
\label{ch:es-donde-dejar-el-coche}
\hlchaptermodality{648}

\begin{chapteropening}
\textbf{By the end of this chapter:} \emph{I can find a parking space, pay at the meter, keep to my lane and talk about trailers and towing.}

Complete the last lesson of chapter 648: I can find a parking space, pay at the meter, keep to my lane and talk about trailers and towing.
\end{chapteropening}

\section[el parquímetro]{el parquímetro — the parking meter}

\label{lesson:ES-C648-parquimetro}



\begin{quote}
Before the new one: say the Spanish for the pedestrian, then the Spanish for the roundabout.
\end{quote}

\subsection*{You'll want to know: el parquímetro}
\textbf{el parquímetro} — \textquotedblleft{}the parking meter\textquotedblright{}.

\begin{quote}
Pon dinero en el parquímetro.
\end{quote}

\begin{quote}
\emph{Put money in the parking meter.}
\end{quote}

You pay at the machine and leave the ticket in the car.

\begin{cousinweb}
\emph{parquímetro} is built on \emph{parking}, borrowed from English (EtymDB), with \emph{-metro}, from Greek \emph{métron} \textquotedblleft{}measure\textquotedblright{}.
\end{cousinweb}

\begin{grammarlens}[title={What you've built}]
One more word for telling, explaining and describing at B1.
\end{grammarlens}

\subsection*{Your turn}
\begin{itemize}
\raggedright
  \item \emph{Say it:} \emph{el parquímetro}
  \item \emph{Say it:} \emph{el parquímetro}, once more
  \item \emph{Say it:} a sentence of your own with \emph{el parquímetro}
  \item \emph{Recall:} say the Spanish for the pedestrian, then the Spanish for the roundabout, then say \emph{el parquímetro} again
\end{itemize}

\begin{tcolorbox}[breakable,colback=teal!4,colframe=teal!35!black,title={Before you move on}]
What does el parquímetro mean? (\textquotedblleft{}The parking meter\textquotedblright{}.) Say it once more.
\end{tcolorbox}
\section[el estacionamiento]{el estacionamiento — the car park, the parking lot; parking}

\label{lesson:ES-C648-estacionamiento}



\begin{quote}
Before the new one: say the Spanish for the roundabout, then the Spanish for the parking meter.
\end{quote}

\subsection*{You'll want to know: el estacionamiento}
\textbf{el estacionamiento} — \textquotedblleft{}the car park, the parking lot; parking\textquotedblright{}.

\begin{quote}
El estacionamiento del hotel es gratis.
\end{quote}

\begin{quote}
\emph{The hotel car park is free.}
\end{quote}

Spain says \emph{el aparcamiento} or \emph{el parking}.

\begin{cousinweb}
\emph{estacionamiento} is built on \emph{estacionar} \textquotedblleft{}to park\textquotedblright{}, which comes from \emph{estación}, Latin \emph{statiō} \textquotedblleft{}a standing still\textquotedblright{} (Lewis \& Short).
\end{cousinweb}

\begin{grammarlens}[title={What you've built}]
One more word for telling, explaining and describing at B1.
\end{grammarlens}

\subsection*{Your turn}
\begin{itemize}
\raggedright
  \item \emph{Say it:} \emph{el estacionamiento}
  \item \emph{Say it:} \emph{el estacionamiento}, once more
  \item \emph{Say it:} a sentence of your own with \emph{el estacionamiento}
  \item \emph{Recall:} say the Spanish for the roundabout, then the Spanish for the parking meter, then say \emph{el estacionamiento} again
\end{itemize}

\begin{tcolorbox}[breakable,colback=teal!4,colframe=teal!35!black,title={Before you move on}]
What does el estacionamiento mean? (\textquotedblleft{}The car park, the parking lot; parking\textquotedblright{}.) Say it once more.
\end{tcolorbox}
\section[el remolque]{el remolque — the trailer; towing}

\label{lesson:ES-C648-remolque}



\begin{quote}
Before the new one: say the Spanish for the parking meter, then the Spanish for the car park.
\end{quote}

\subsection*{You'll want to know: el remolque}
\textbf{el remolque} — \textquotedblleft{}the trailer; towing\textquotedblright{}.

\begin{quote}
Llevamos el barco en un remolque.
\end{quote}

\begin{quote}
\emph{We are taking the boat on a trailer.}
\end{quote}

For a boat, a horse or a small caravan.

\begin{cousinweb}
\emph{remolque} is built on \emph{remolcar} \textquotedblleft{}to tow\textquotedblright{}, from a Vulgar Latin verb (EtymDB) built on Latin \emph{remulcum} \textquotedblleft{}a tow-rope\textquotedblright{} (Lewis \& Short).
\end{cousinweb}

\begin{grammarlens}[title={What you've built}]
One more word for telling, explaining and describing at B1.
\end{grammarlens}

\subsection*{Your turn}
\begin{itemize}
\raggedright
  \item \emph{Say it:} \emph{el remolque}
  \item \emph{Say it:} \emph{el remolque}, once more
  \item \emph{Say it:} a sentence of your own with \emph{el remolque}
  \item \emph{Recall:} say the Spanish for the parking meter, then the Spanish for the car park, then say \emph{el remolque} again
\end{itemize}

\begin{tcolorbox}[breakable,colback=teal!4,colframe=teal!35!black,title={Before you move on}]
What does el remolque mean? (\textquotedblleft{}The trailer; towing\textquotedblright{}.) Say it once more.
\end{tcolorbox}
\section[el carril]{el carril — the lane; the rail}

\label{lesson:ES-C648-carril}



\begin{quote}
Before the new one: say the Spanish for the car park, then the Spanish for the trailer.
\end{quote}

\subsection*{You'll want to know: el carril}
\textbf{el carril} — \textquotedblleft{}the lane; the rail\textquotedblright{}.

\begin{quote}
El carril de la derecha es para autobuses.
\end{quote}

\begin{quote}
\emph{The right-hand lane is for buses.}
\end{quote}

\emph{El carril bici} — the cycle lane.

\begin{cousinweb}
\emph{carril} is usually traced to \emph{carro} \textquotedblleft{}cart\textquotedblright{}, from Latin \emph{carrus}; it first named the rut that a cart wheel leaves on a road.
\end{cousinweb}

\begin{grammarlens}[title={What you've built}]
One more word for telling, explaining and describing at B1.
\end{grammarlens}

\subsection*{Your turn}
\begin{itemize}
\raggedright
  \item \emph{Say it:} \emph{el carril}
  \item \emph{Say it:} \emph{el carril}, once more
  \item \emph{Say it:} a sentence of your own with \emph{el carril}
  \item \emph{Recall:} say the Spanish for the car park, then the Spanish for the trailer, then say \emph{el carril} again
\end{itemize}

\begin{tcolorbox}[breakable,colback=teal!4,colframe=teal!35!black,title={Before you move on}]
What does el carril mean? (\textquotedblleft{}The lane; the rail\textquotedblright{}.) Say it once more.
\end{tcolorbox}
\section[remolcar]{remolcar — to tow}

\label{lesson:ES-C648-remolcar}



\begin{quote}
Before the new one: say the Spanish for the trailer, then the Spanish for the lane.
\end{quote}

\subsection*{You'll want to know: remolcar}
\textbf{remolcar} — \textquotedblleft{}to tow\textquotedblright{}.

\begin{quote}
La grúa remolca el coche hasta el taller.
\end{quote}

\begin{quote}
\emph{The tow truck tows the car to the garage.}
\end{quote}

The c becomes qu before e: \emph{remolqué}. \emph{El remolque} is a trailer.

\begin{cousinweb}
\emph{remolcar} comes from a Vulgar Latin verb (EtymDB) built on Latin \emph{remulcum} \textquotedblleft{}a tow-rope\textquotedblright{} (Lewis \& Short), a word of seafaring that the Romans took from Greek.
\end{cousinweb}

\begin{grammarlens}[title={What you've built}]
One more word for telling, explaining and describing at B1.
\end{grammarlens}

\subsection*{Your turn}
\begin{itemize}
\raggedright
  \item \emph{Say it:} \emph{remolcar}
  \item \emph{Say it:} \emph{remolcar}, once more
  \item \emph{Say it:} a sentence of your own with \emph{remolcar}
  \item \emph{Recall:} say the Spanish for the trailer, then the Spanish for the lane, then say \emph{remolcar} again
\end{itemize}

\begin{tcolorbox}[breakable,colback=teal!4,colframe=teal!35!black,title={Before you move on}]
What does remolcar mean? (\textquotedblleft{}To tow\textquotedblright{}.) Say it once more.
\end{tcolorbox}
